\documentclass[10pt]{article}

\usepackage[a4paper,left=1.6cm,right=1.6cm,top=1.6cm,bottom=1.6cm]{geometry}
\usepackage{enumitem}
\usepackage{titlesec}
\usepackage{xcolor}
\usepackage{hyperref}
\usepackage{fontawesome5}
\usepackage{array}

\definecolor{accent}{HTML}{0E7490}
\hypersetup{colorlinks=true, urlcolor=accent, linkcolor=accent}

\titleformat{\section}{\large\bfseries\color{accent}}{}{0em}{}[\vspace{-0.6em}\rule{\textwidth}{0.6pt}]
\titlespacing{\section}{0pt}{0.9em}{0.5em}
\setlist[itemize]{leftmargin=1.2em,itemsep=0.15em,topsep=0.2em,parsep=0pt}

\newcommand{\entry}[4]{%
  \textbf{#1}\hfill #2\\
  \textit{#3}\hfill #4\par\vspace{0.2em}%
}

\pagestyle{empty}

\begin{document}

\begin{center}
  {\Huge\bfseries 张\ 三}\\[0.3em]
  {\color{accent}\small \faPhone\ 138-0000-0000 \quad \faEnvelope\ zhangsan@example.com \quad \faGithub\ github.com/zhangsan \quad 北京}
\end{center}

\section{教育经历}
\entry{某某大学}{2020.09 -- 2024.06}{计算机科学与技术，本科}{GPA 3.8/4.0}
\begin{itemize}
  \item 主修课程：数据结构、操作系统、机器学习、编译原理。
  \item 获得校级一等奖学金两次。
\end{itemize}

\section{项目经历}
\entry{在线 LaTeX 协作编辑器}{2024.03 -- 2024.09}{核心开发者}{React / Node.js}
\begin{itemize}
  \item 基于 WebSocket 与 CRDT 实现多人实时协同编辑，冲突率降至 0。
  \item 设计增量编译调度，将平均编译等待时间缩短 40\%。
\end{itemize}

\entry{分布式日志采集系统}{2023.10 -- 2024.02}{后端开发}{Go / Kafka}
\begin{itemize}
  \item 支撑每日 2 亿条日志的采集与检索，P99 延迟低于 200ms。
\end{itemize}

\section{专业技能}
\begin{itemize}
  \item \textbf{编程语言：} C/C++、Python、JavaScript、Go、\LaTeX
  \item \textbf{工具与框架：} Git、Docker、React、PyTorch
  \item \textbf{语言：} 英语 CET-6，可流畅阅读英文技术文档
\end{itemize}

\end{document}
